\section{A finite formula for strip crossing mass}\label{c6:strip-transfer}

The local matrix is a specialization of the dilute-loop integrable algebra. Grimm and Pearce~\cite[Sections~4.1 and~4.6]{GrimmPearce1993} construct dilute Yang--Baxter operators by braid--monoid methods. They explain that local gauge transformations preserve doubly periodic partition functions, while other quantities can depend on the gauge. We fix the gauge and boundary vectors explicitly: every identity used to evaluate the bridge mass is verified below.


Use the boundary midpoints and visited-center length convention of Section~\ref{sec:introduction}. For the transfer calculation set
\[
 L=\pi/8,\quad S=\sin(2L)=1/\sqrt{2},\quad
 b=3L,\quad \gamma=b/2,\quad
 \kappa=(2\cos L)^{-1}.
\]
Take lattice axes at $60^\circ$, grouped into unit lozenges (edges called w,s,e,n in the two lattice coordinates, with w,s meeting at an acute corner). Consider the infinite strip with first coordinate between 0 and $N$. Its two boundary cuts are at Euclidean distance $Nd_0$, so after a lattice rotation its crossing mass is the quantity $B_h$ of Section~\ref{sec:introduction} with $h=N$. Let $B_N$ be the sum of weights from a fixed edge midpoint on the right side to all possible edge midpoints on the left side, through the interior. The following finite-dimensional formula for $B_N$ is useful. (Brackets below mean bilinear integration, not conjugating coefficients.)

The function $C_m(x)=\cosh(m\sqrt{x})$ is entire in $x$, as its power series shows; in particular its values at negative $x$ require no choice of square root. Set $D(x)=2C_L(x)-1$, $a_*=D(-4)=2\cos(2L)-1$. Use the positive measure $\mu$ parametrized by
\[
 d\mu(x)=\frac{zD(x)}{\sinh(4Lz)}\,dz,\qquad x=z^2,\quad z\in (0,\infty).
\]
Let
\[
 M=\operatorname{span}\{1,x,\ldots,x^{N-1};\ \partial_m^j C_m(x)|_{m=L},\quad j=0,\ldots,N-1\},\qquad
 g(x)=\sinh(L\sqrt{x})/\sqrt{x}.
\]
\begin{theorem}[Finite strip functional]\label{thm:strip-functional}
For every integer $N\ge1$, there exist $\rho\in\mathbf C$ and a functional of the form
\[
 \mathcal L(H)=4a_* H(-4)+\rho H(0)+\int H f\,d\mu,\qquad f\in M+\mathbf C g
\]
defined for functions $H$ with specified values at $-4$ and $0$ whenever the displayed integral converges absolutely, such that
\begin{equation}
 \begin{split}
 \mathcal L(M)=0,\qquad \mathcal L(C_{L/2})=0,\qquad
 \mathcal L(C_{3L/2}-C_{L/2}/D)=0,\qquad 2 B_N=\mathcal L(1/D).
 \end{split}                                                     \label{eq:legacy-P-1}
\end{equation}
\end{theorem}
We first construct a commuting double-row family from local bulk and boundary identities. Its mixed boundary coefficient will give a scalar function whose value at the physical weights is $2B_N$. The scalar crossing lemma in Section~\ref{sec:strip-scalar} records the poles, growth, and fusion zeros of this function; these are the interpolation data needed to prove Theorem~\ref{thm:strip-functional}.

\subsection{Local matrices and fusion}\label{sec:strip-local-identities}
Let the three basis symbols be $+,0,-$ identified with spins $1,0,-1$. Write $s=\exp(2iL)$. For spectral parameter $u$ define normalized coefficients
\[
 \begin{aligned}
 (a,h,v,c,d)(u)&=\frac{1}{Q(u)}
 \big(S\sin(b-u),\ S\sin u,\ \sin(b-u)\sin u,\;\\
 &\hspace{3em}\sin(b-u)\sin(2L-u),\ \sin u\sin(u-L)\big),\\
 Q(u)&=\sin(2L+u)\sin(b+u).
 \end{aligned}
\]
The matrix $R(u)$ on an ordered pair of spin spaces acts by $c$ on $++$ and $--$; on each of $\operatorname{span}(+0,0+)$ and $\operatorname{span}(-0,0-)$ by $\left(\begin{smallmatrix}v&a\\a&v\end{smallmatrix}\right)$; and on $(+-,00,-+)$ by
\[
 \begin{pmatrix}
 d&h/s&c+d/s^2\\ h s&1&h/s\\ c+d s^2&h s&d
 \end{pmatrix}.
\]
Use boundary variables $e$ with $e^2=0$, commuting. Define
\[
 \begin{gathered}
 k(t)=\frac{\sin(\gamma-t)}{\sin(\gamma+t)},\quad
 l(t)=\frac{\sin(2t)}{\sin(\gamma+t)},\\
 D_e(t)=\operatorname{diag}(k,1,k)+ e\,l(E_{+,0}+s E_{0,-}),\\
 K_e(t)_{p,-m}=s^m (D_e(t))_{p m},
 \end{gathered}
\]
where $E_{ij}$ are matrix units and $K_e$ is used as a two-spin vector. On a row of $N$ (site) spins with auxiliary spin $\alpha$, rapidities $\xi_i$, put
\[
 U_\alpha(t)=R_{\alpha 1}(t+\xi_1)\cdots R_{\alpha N}(t+\xi_N).
\]
Subscripts indicate the ordered spaces acted upon; other factors are spectators. Contract two auxiliary rows in the double row operator ($r=b-t$)
\[
 T(t) = K_{e_\ell}(r)_{AB}^{T} U_A(t) U_B(r) K_{e_r}(t)_{AB}.
\]
Here $e_\ell^2=e_r^2=0$ separately, and superscript $T$ is ordinary transpose. The nilpotent variables retain at most one boundary insertion on each side. Figure~\ref{fig:strip-double-row} shows the order of the two auxiliary rows and the site spaces.

\begin{figure}[t]
\centering
\begin{tikzpicture}[scale=.88,>=Latex,font=\small]
  \draw[thick] (0,.65) -- (5,.65);
  \draw[thick] (0,1.85) -- (5,1.85);
  \draw[->] (4.65,.65) -- (4.25,.65);
  \draw[->] (4.65,1.85) -- (4.25,1.85);
  \foreach \x/\lab in {1/1,2.15/2,4/N}{
    \draw[thick,->] (\x,-.05) -- (\x,2.55);
    \draw[fill=white] (\x-.12,.53) rectangle (\x+.12,.77);
    \draw[fill=white] (\x-.12,1.73) rectangle (\x+.12,1.97);
    \node[below] at (\x,-.05) {$\lab$};
  }
  \node at (3.05,1.25) {$\cdots$};
  \draw[thick] (0,.65) .. controls (-.55,.65) and (-.55,1.85) .. (0,1.85);
  \draw[thick] (5,.65) .. controls (5.55,.65) and (5.55,1.85) .. (5,1.85);
  \node[anchor=east,align=right] at (-.5,1.25) {$K_{e_\ell}(r)^T$};
  \node[anchor=west] at (5.5,1.25) {$K_{e_r}(t)$};
  \node[above] at (3.25,1.88) {$A:\ t$};
  \node[below] at (2.5,-.85) {$B:\ r=b-t$};
  \node[above] at (2.5,2.7) {Double-row contraction};
  \begin{scope}[shift={(8.3,.35)},scale=1.15]
    \coordinate (O) at (0,0);
    \coordinate (V) at (1.5,0);
    \coordinate (U) at (.75,1.299);
    \coordinate (UV) at (2.25,1.299);
    \coordinate (C1) at (.75,.433);
    \coordinate (C2) at (1.5,.866);
    \draw[gray] (O)--(V)--(UV)--(U)--cycle;
    \draw[gray] (U)--(V);
    \draw[thick] (.375,.6495)--(C1)--(.75,0);
    \draw[thick] (C1)--(C2);
    \draw[thick] (1.5,1.299)--(C2)--(1.875,.6495);
    \fill (C1) circle (1.6pt);
    \fill (C2) circle (1.6pt);
    \node[left] at (.375,.6495) {$w$};
    \node[below] at (.75,0) {$s$};
    \node[above] at (1.5,1.299) {$n$};
    \node[right] at (1.875,.6495) {$e$};
    \node[above] at (1.125,1.6) {Physical lozenge};
  \end{scope}
\end{tikzpicture}
\caption{The horizontal and vertical lines denote spin spaces, and the boundary curves denote the displayed tensors. The auxiliary row $B$ acts first on the site spins, then row $A$; products act from right to left. The two auxiliary ends are contracted with the displayed boundary tensors. This is a tensor schematic, whereas the inset shows the actual pair of equilateral triangles at the physical specialization. The black segments in the inset are the available honeycomb links, not one path configuration. The two dots are the centers counted by the path length: an $s$--$w$ or $e$--$n$ arc visits one center, and a single arc between the two pairs visits both.}
\label{fig:strip-double-row}
\end{figure}


We record local identities and their algebraic verification. Equalities are meromorphic (used first away from singularities). Let $P$ swap a pair of spin spaces, $F$ flip spins in a space, $J=\operatorname{diag}(-1,1,-1)$. Define
\[
 \begin{gathered}
 \Omega=\sum_p s^{-p}|p,-p\rangle,\qquad
 W|\pm\rangle=\kappa(|\pm,0\rangle+|0,\pm\rangle),\\
 W|0\rangle=|00\rangle+\kappa\sum_{p=\pm1}s^{-p}|p,-p\rangle,\quad Z=W^*=\bar W^T.
 \end{gathered}
\]
The image slots of $W$ will be indicated by a pair of labels (ordered); its input slot can then be regarded as a fused spin. We have
\begin{equation}
\begin{gathered}
 R_{21}(u)=R_{12}^T(u),\qquad R_{12}(u)R_{12}^T(-u)=I,\\
 R(u+\pi)=J_1 R(u)J_1,\qquad
 R(b-u)_{ij,kl}=s^{k-i}R(u)_{-i,l;-k,j},\\
 R(0)=P,\qquad R(b)=\Omega\bar\Omega^T,\qquad R(2L)=W Z.
\end{gathered}                                                   \label{eq:legacy-P-2}
\end{equation}
Indeed $Q(b-u)=Q(u)$, swapping $a,h$ and $c,d$ at reflected argument, which gives crossing by the displayed entries. Unitarity on the $\pm2,\pm1$ blocks uses $c(u)c(-u)=1$, $v(u)/a(u)=\sin u/S$ and $a(u)a(-u)=S^2/(S^2-\sin^2 u)$. On the zero block use additionally (prime denotes argument $-u$ here) $h+c h'=0$, $h h'+c d'+c' d=0$, by the sine products, and $s^2+s^{-2}=0$. Values in \eqref{eq:legacy-P-2} follow by substitution. Consequently
\[
 R_{23}(u-b)R_{13}(u)\Omega_{12}=\Omega_{12} I_3,
\]
where $I_3$ denotes the identity on slot 3 (by crossing, the output $(p,q,k)$ at input $j$ is $s^{-p}(R(u-b)R(b-u)^T)_{qk;-p,j}$).

The basic fusion identity and its adjoint are
\begin{equation}
 R_{23}(u-L)R_{13}(u+L)W_{12}=W_{12} R_{f3}(u),\qquad
 Z_{12} R_{13}(u+L)R_{23}(u-L)=R_{f3}(u) Z_{12},                \label{eq:legacy-P-3}
\end{equation}
where $f$ labels the fused space. Here and below adjoints can be taken at nonsingular real parameters first ($R$ Hermitian). We give details of the elementary check to fix normalizations. Multiply the first formula by $R_{32}(L-u)$ using \eqref{eq:legacy-P-2}, and write a subscript $+$ for argument $u+L$, primes for $L-u$, no decorations for $u$. Expanding on inputs $(f,3)$ gives the following sufficient equalities (spin-flipped inputs conjugate the same ones):
\[
\begin{array}{l|l}
 ++&cv'=c_+,\ cc'=v_+,\ ca'=a_+\\
 +0&v+a h'=v_+,\ vv'+aa'/\kappa=1,\ va'+av'/\kappa=a_+,\;\\
 &v h'+a d'=0,\ ac'=h_+\\
 0+&\kappa(a+v h')=a_+,\ av'+va'/\kappa=h_+,\ \kappa aa'+vv'=v_+,\;\\
 &a h'+v d'=c_+,\ vc'=d_+\\
 00&\kappa(a'+v'h)=h_+,\ v'+a'h=v_+,\ h'+\kappa h c'=\kappa a_+\\
 +-&h+\kappa c h'=\kappa h_+,\ hh'+\kappa(c d'+c'd)=0,\\
 &d v'+h a'=d_+,\;d a'+h v'=0,\text{ and the }++\text{ identities}.
\end{array}
\]
To verify, put $t_j=\sin(u+jL)$; the numerator rows and denominators are
\[
\begin{array}{c|rrrrr|l}
 &a&h&v&c&d&Q\\
 &-S t_{-3}&S t_0&-t_{-3}t_0&t_{-3}t_{-2}&t_0t_{-1}&t_2t_3\\
 '&S t_2&-S t_{-1}&-t_{-1}t_2&t_2t_1&t_{-1}t_0&-t_{-3}t_4\\
 +&-S t_{-2}&S t_1&-t_1t_{-2}&t_{-2}t_{-1}&t_1t_0&t_3t_4 .
\end{array}
\]
Thus $QQ'/Q_+=-t_{-3}t_2$. The $++$ row and $v h'+a d'=0$, $ac'=h_+$, $vc'=d_+$ follow immediately; for the other identities in rows $+0,0+$ use $t_{-2}t_2=t_0^2-S^2$, $t_{j-1}+t_{j+1}=t_j/\kappa$,
$t_{-3}t_4=t_0t_1-S^2/\kappa$, $t_3t_4=S^2/\kappa-t_0t_{-1}$, $t_0t_{-1}-t_{-2}t_1=\kappa S^2$. For the last two rows one also has
$t_2t_1-t_{-2}t_{-1}=t_0t_4/\kappa$,
$t_{-3}t_{-2}+t_2t_1=S^2/\kappa$, and can interchange primed and unprimed quantities (swaps $h_+,a_+$ and $d_+,c_+$).

These imply Yang--Baxter:
\[
 R_{12}(u-v) R_{13}(u) R_{23}(v)=R_{23}(v) R_{13}(u) R_{12}(u-v).
\]
Indeed at $u-v=0$ use the permutation; at $u-v=2L$ use $R(2L)=W Z$ and \eqref{eq:legacy-P-3}; at $u-v=b$ use \eqref{eq:legacy-P-2} and the analogous rank-one identity shown after \eqref{eq:legacy-P-2}, and its adjoint. At $u=0,2L$ reduce to the first two cases by multiplying on both sides by $R_{23}(v)^{-1}=R_{32}(-v)$ and permuting the labels (new order $1,3,2$). For generic fixed $v$ these give five values modulo $\pi$; clearing denominators, each entry has trigonometric degree $\le4$ in $u$ and is either periodic or antiperiodic by $\pi$ by \eqref{eq:legacy-P-2}, so vanishes identically.

\paragraph{Boundary fusion and reflection.}
The bulk identities are now available. The boundary identities needed to close and exchange rows (using a common variable $e$ for each equality) are, with $r=b-t$,
\[
 D_e(t)D_e(-t)=I,\qquad R(t-r) K_e(t)=P K_e(r),
\]
\begin{equation}
 \begin{split}
 D_{e,2}(t-L)R(2t)D_{e,1}(t+L)\bar W&=W D_e(t),\\
 W^T D_{e,2}(t+L)R(2t)D_{e,1}(t-L)&=D_e(t)\bar W^T .
 \end{split}                                                       \label{eq:legacy-P-4}
\end{equation}
In $D_{e,j}$ the second subscript just labels a slot. Unitarity of $D_e$ follows from $k(-t)=1/k(t)$, $l(-t)=-l(t)/k(t)$. To swap $K$, on total spin 0 and 1 the equalities reduce to $c(t-r) k(t)+h(t-r)=k(r)$ and $(v+a)(t-r)l(t)=l(r)$. Writing $y=t-\gamma$, these follow by
\[
 \sin(L-y)\sin(b-2y)+\sin(L+y)\sin(b+2y)=S\cos y,\qquad
 \cos(L\pm y)=\sin(b\mp y)
\]
on inserting the weights. For \eqref{eq:legacy-P-4}, the second identity follows from the first at real $t$ by conjugation, flipping all spins and swapping slots 1 and 2, then transposing ($J$ is not needed: $F\bar D_e F=D^T_{e/s}$ using formal real $e$). For the first, write subscripts $+,-$ on $k,l$ for $t\pm L$ and evaluate $a,h,v,c,d$ at $2t$. On expanding the three columns, the constant and linear terms in $e$ reduce respectively to
\[
 \begin{split}
 v k_+ +a=k,\qquad k_-(a k_+ +v)=k,\qquad k_-(h+\kappa c k_+)=\kappa;\\
 l_+ k_- c+k l_-=0,\qquad l_+(v+\kappa a)+\kappa l_- /k_-=\kappa l,\qquad
 l_+ k_-(a+\kappa v)+l_-=\kappa l,\\
 l_+ k_-(h+d)=l,\qquad \kappa\big((1+h)l_++k l_-/k_-\big)=l .
 \end{split}
\]
For example put $H_\pm=\sin(\gamma\pm t)$, $G_\pm=\cos(\gamma\pm t)$, $F_\pm=\sin(L/2\pm t)$, $P_\pm=\sin(2L\pm2t)$, $T_0=\sin(2t)$. Then $k_+=F_-/G_-$, $k_-=G_+/F_+$, $l_+=P_+/G_-$, $l_-=-P_-/F_+$; substitution gives the equalities by the sine products
\[
\begin{gathered}
 P_+ G_+-S G_-=T_0 F_-,\qquad P_+F_+-S F_-=T_0 G_-,\qquad
 2(H_+ F_+ P_+ -H_-F_-P_-)= S T_0/\kappa,\\
 P_- H_+-S H_-=T_0 F_-,\quad \kappa H_- -F_-=\kappa F_+,\quad H_-=\kappa(F_-+G_+),\quad
 S+\sin(2t-L)=2F_+G_-,\\
 P_+ H_+ G_+-P_- H_- G_-=\cos L\,T_0\cos(2t).
\end{gathered}
\]

We will use the reflection relation
\begin{equation}
 R(t-y) D_{e,1}(t) R(t+y)^T D_{e,2}(y)
 =D_{e,2}(y)R(t+y)D_{e,1}(t)R(t-y)^T.                            \label{eq:legacy-P-5}
\end{equation}
Here is a polynomial check using \eqref{eq:legacy-P-4}. Fix generic $y$. At $t=y,-y$ use $R(0)=P$ and $D_e(y)\allowbreak D_e(-y)=I$; at $t=0,4L$ use $D_e=I,J$ respectively, and \eqref{eq:legacy-P-2}. At $t-y=2L$ the right side is $W D_e((t+y)/2)W^T$ by \eqref{eq:legacy-P-4} and the left is the same (swap 1,2 in the second relation of \eqref{eq:legacy-P-4}, using $PW=\bar W$). This also gives \eqref{eq:legacy-P-5} at $t+y=2L$ using $y\mapsto-y$ and $D_e(y)D_e(-y)=I$. At $t-y=b$ the right side is $\Omega \Omega^T$, since
\[
 D_{e,2}(y) R(t+y)D_{e,1}(t)\bar\Omega
 =D_{e,2}(y)R(t+y)K_e(t)
 =P D_{e,1}(y)K_e(-y)=\Omega .
\]
This last equality also holds with both $D$'s transposed (flip both spins and conjugate, rescaling $e$), giving the same left side of \eqref{eq:legacy-P-5}. After clearing denominators the difference in \eqref{eq:legacy-P-5} has degree $\le 6$ in $t$, entries periodic or antiperiodic by $\pi$ (after clearing, $D_{e,1}$ transforms by conjugation by $J_1$ and a sign $-$); hence the seven checks suffice.

\subsection{Commuting and fusing double rows}\label{sec:strip-double-row-identities}
We now derive identities for $T$ by row exchange. Yang--Baxter gives
$R_{ij}(u-v) U_i(u)U_j(v)=U_j(v)U_i(u) R_{ij}(u-v)$. In particular swapping the two rows of $T$ gives $T(t)=T(r)$: the input swap uses $R(t-r)K_e(t)=P K_e(r)$ and the output swap its $r,t$ version via unitarity. Also $T(t+\pi)=T(t)$ since $K_e(t+\pi)=J_1 J_2 K_e(t)$.

For $T(t)T(y)$, take labels $A,B,C,D'$ of parameters $t,r,y,b-y$ respectively (below write $D'$ for a label). To interchange the row pairs, exchange $BC,AC,BD',AD'$ in that order, acting on the right boundary vector $V=K_e(t)_{AB}K_e(y)_{CD'}$ ($e=e_r$) by the corresponding $R$ product. Invariance of the vector reduces by \eqref{eq:legacy-P-2} to
\[
 R_{AC}(t-y)R_{BC}(b-t-y)V
 =R_{BD'}^T(t-y)R_{AD'}^T(b-t-y)V .
\]
This is \eqref{eq:legacy-P-5} by crossing. Indeed identify a vector $X$ with the matrix with entries $s^{-m-n}X_{i,-m,j,-n}$ in row $i,j$ and column $m,n$. Under this identification $R_{BC}(b-t-y)V$ is $D_{e,1}(t) R(t+y)^T D_{e,2}(y)$ and $R_{AD'}^T(b-t-y)V$ is $D_{e,2}(y) R(t+y) D_{e,1}(t)$. The remaining $R$'s act on the left and right respectively as in \eqref{eq:legacy-P-5}. On the left boundary (covector written as a transposed column), the inverse-transposed exchanges appear, exactly the analogous computation with $t,r$ exchanged and $y,b-y$ exchanged, $e=e_\ell$. Thus all $T$ commute. Further $T(0)=T(4L)=I$: at $0$ the input is $K_e(0)=\bar\Omega=P\Omega$ and the two rows act trivially on it by rank-one crossing after \eqref{eq:legacy-P-2}. The output pairing with this vector is 1. At $4L$, shift the second parameter $r$ by $\pi$ (conjugate that slot by $J$); the input becomes again $\bar\Omega$, the parameter difference again $b$ (higher parameter second), and the output pairing is 1.

For a product $T(t-L)T(t+L)$ use labels $B,C,A,D'$ with parameters $t-L,r+L,t+L,r-L$ in that order. Swap middle rows by $R_{CA}(r-t)$. We obtain product order $B A C D'$ between boundary column vectors $G_r,G_\ell$ (output transposed):
\[
 \begin{split}
 G_r &= R_{CA}(b-2t) K_{e_r}(t-L)_{BC} K_{e_r}(t+L)_{AD'},\\
 G_\ell &= R_{CA}(b-2r) K_{e_\ell}(r+L)_{BC} K_{e_\ell}(r-L)_{AD'} .
 \end{split}
\]
With fused slots $f,g'$ for $AB,C D'$ respectively, \eqref{eq:legacy-P-4} gives
\[
 (1\otimes Z_{C D'})G_r=(W_{AB}\otimes1)K_{e_r}(t)_{f g'},\qquad
 (W_{AB}^T\otimes1)G_\ell=(1\otimes Z_{C D'}^T)K_{e_\ell}(r)_{f g'} .
\]
Indeed the matrix given by $s^{-m-n} (G_r)_{p q;-m,-n}$ (slots indexed $AB;C D'$, columns ordered $n,m$) by crossing is $D_{e_r,2}(t-L) R(2t) D_{e_r,1}(t+L)$, and similarly for $G_\ell$ with $r,+L,-L$. Also $Z_{-j;-m,-n}=\bar W_{n,m;j}$ and it is zero unless $j=m+n$.

By \eqref{eq:legacy-P-3}, $U_B U_A$ preserves the $W_{AB}$ subspace acting by $U_f(t)$, and $U_C U_{D'}$ preserves $\ker Z_{C D'}$ acting on the quotient by $U_{g'}(r)$. Moreover:

\begin{itemize}

\item 

At $t=L-\xi_i$ the local product $R_{Bi}(0) R_{Ai}(2L)$ has range in $W_{AB}$, hence so does $U_B U_A$.

\item 

At $t=L+\xi_i$ the local product $R_{Ci}(b) R_{D'i}(L)$ kills $\ker Z_{C D'}$, hence so does $U_C U_{D'}$. Indeed its adjoint has range in the image of $W_{C D'}$: $R_{D'i}(L)$ maps $\Omega_{Ci}$ into this image by crossing at $u=2L$, $R(2L)=WZ$ (output $c,d,j$ for input $v$ on $D'$ has coefficient $s^{-c}R(L)_{dj;v,-c}=s^j R(2L)_{cd;-j,v}$).

\end{itemize}

Consequently at either of those values the contracted product reduces to $T(t)$. Indeed the part of $G_r$ in arbitrary slots $AB$ tensor $\ker Z_{CD'}$ is then harmless (either killed, or mapped to $W_{AB}\otimes\ker Z_{CD'}$ annihilated by the output); the rest can be lifted through $W_{AB}$ from $K_{e_r}(t)_{f g'}$, and \eqref{eq:legacy-P-3} and the formula for $G_\ell$ give precisely $T(t)$. We have shown, first for generic rapidities,
\begin{equation}
 T(t-L)T(t+L)=T(t)\quad\text{at }t=L\pm\xi_i.                     \label{eq:legacy-P-6}
\end{equation}

\subsection{The scalar crossing function}\label{sec:strip-scalar}

The row identities now have to be connected to the positive path sum. The next statement collects the scalar data that will enter the finite functional.

\begin{auxlemma}[Scalar crossing function]
At the homogeneous rapidities $\xi_1=\cdots=\xi_N=-L/2$, the double-row family determines an even, $\pi$-periodic meromorphic function $p$ with $p(0)=2B_N$. It is rational in $X=\cos(2u)$, with polynomial part $4X+\rho$ for a constant $\rho$, and denominator dividing
\[
 (X+1)^N(X+\cos 2L)^{N+1}.
\]
Moreover $p(3L/2)=p(5L/2)=0$, and
\[
 p(u-L)+p(u+L)-p(u)
\]
has a zero of order at least $2N$ at $u=0$ and of order at least $N$ at $u=L$.
\end{auxlemma}
\begin{proof}
We construct the scalar function through the physical path expansion and then verify its rational form and zero multiplicities.

\paragraph{Loop cancellation.}
The finite matrix identities alone do not identify a walk partition function. To extract the relevant eigenvalue of the commuting family, specialize first to $e_\ell=e_r=0$. Matrix products have a planar oriented-path expansion on the square tensor diagram of Figure~\ref{fig:strip-double-row}. Positive horizontal arrows are right-to-left and positive vertical arrows upward. In a tile, allowed arc configurations, with orientation ignored, are empty (weight 1), a single e-n or s-w ($a$), a single e-s or n-w ($h$), a single straight across ($v$), or both e-n, s-w ($c$) or both e-s, n-w ($d$). The entries of $R$ give these weights times phases, summed over choices: all arc phases are 1 except for e-to-s, whose phase is $s$, and n-to-w, whose phase is $s^{-1}$ (reverse arrows conjugate the phases). At the left and right boundaries $K_0$ permits empty or a U-turn between the double row's two terminals, the latter with weight $k$ and phase $s$ if the upper arrow is positive ($s^{-1}$ otherwise). These phase assignments have product $s^{\pm2}$ on any simple closed path, and product 1 on a simple bridge from a right endpoint to a left endpoint in a rectangle, with only the hard walls (no U-turns). For instance give each arc phase $\exp(i\,\Delta/4)$ where $\Delta$ is its rotation of tangent (starting and ending orthogonally to tile edges), plus $\pi/2$ at terminal vertical arrows, minus $\pi/2$ at initial vertical arrows. These corrections telescope; take quarter circles and straight lines and smooth boundary turns as needed. Bridges have total turn 0 by simplicity (close using the boundary). Thus, summing orientations, closed loops cancel. This holds for arbitrary spectral parameters at zero boundary variables. In particular every product of double rows at those variables has empty-site-to-empty-site entry 1.

At $t=\gamma$ and $\xi_i=-L/2$, the boundary weights $k$ are zero, and the tile weights are evaluated at $L$, giving $a=\kappa$, $h=v=c=\kappa^2$, $d=0$. These are precisely hard-wall honeycomb paths through triangular pairs: s-w and e-n each visit the center of one triangle, distinct when both are used; each of the other single arcs visits both centers. The square drawing is topological; the inset of Figure~\ref{fig:strip-double-row} shows the actual honeycomb geometry.

\paragraph{A signed flux bound.}\label{sec:strip-flux}
We next obtain positive path bounds at these weights. The calculation is the local parafermionic cancellation of \cite[Lemma~1]{DuminilCopinSmirnov2012}, written in our signed exit-flux convention. It uses the physical honeycomb turning angle, rather than the oriented-loop phases above. In any finite lozenge rectangle and with one starting boundary midpoint, the sum of weights of paths to the rest of the boundary is at most $1/\cos(3\pi/8)$. Indeed insert in each such path weight (including in prefixes) the factor $\exp(i\beta W_{\rm turn})$, $\beta=3/8$, $W_{\rm turn}$ total turn in actual honeycomb geometry. At each triangle center, the total signed flux of prefixes into it (ending at an incident midpoint heading toward it, including the initial flux 1 at the source) equals the flux out of it. If unvisited the two continuations sum back to the entering weight by $2\kappa\cos L=1$. The incoming terms where the vertex has been visited cancel in pairs: there is a stem into the vertex followed by the unfinished simple loop back to its third half-edge; reversing the loop, the changes of turn from the stem are $+4\pi/3$ and $-4\pi/3$. Indeed the stem arrives from outside the loop by simplicity and the boundary start; for a completed positively oriented loop the turn at this vertex would be $+\pi/3$, but departure there from the stem instead gives $-\pi/3$. (Weights in the pair have thus opposite phases.) This proves the flux rule, hence total signed exit mass 1. Turns to the boundary are in $[-\pi,\pi]$ by convexity and simplicity (close from the end to the start along the perimeter; besides the arc turn there are two turns by $\pi/2$ with common sign and the boundary direction change). This gives the bound by real parts.

\paragraph{Decay at fixed width.}
With width fixed, path weight from a fixed bottom midpoint to the opposite top boundary of taller and taller lozenge rectangles tends to 0. Indeed to each top exit append in an extra row a deterministic continuation to a lateral wall (one turn s-w and then e-w straight crossings as necessary), staying there in new triangles, at cost bounded below depending only on width. This is injective for each height and sends different heights to different endpoints. Total exit mass to the lateral boundaries in the semi-infinite rectangle is bounded by the preceding flux argument and exhaustion. All bottom arch weights in the semi-infinite rectangle (analogously top arch weights) are summable by the same bound.

\paragraph{The common eigenline.}
Now every entry of $T(\gamma)^m$ in this specialization is an oriented sum of path systems between the specified occupied top/bottom midpoints of the rectangle of height $2m$ (no detached loops). Systems containing at least one bottom-to-top path contribute vanishingly as $m\to\infty$, by bounding by products of the preceding single-path sums; the pure bottom and top arches converge to independent half-strip systems by the summability (phases included). Consequently $T(\gamma)^m$ tends to a rank-one projection with empty-empty entry 1. Writing $\mathbf0$ for the empty site-spin state, we have
\[
 T(\gamma)^m\longrightarrow P_0,\qquad \rank P_0=1,\qquad (P_0)_{\mathbf0,\mathbf0}=1.
\]
Thus $1$ is a simple eigenvalue and the rest of the spectrum lies strictly inside the unit circle.

For rapidities near this specialization we can analytically continue that simple eigenline of $T(\gamma)$ (and its eigenprojection), also over the square-zero boundary variables by expansion. At zero boundary variables the simple eigenvalue is still 1: all empty-empty entries of powers remain 1 by loop cancellation, while the other spectral part decays. The corresponding projection thus has empty-empty entry 1. 

Since all $T(t)$ commute (modulo $e_\ell^2=e_r^2=0$), this is a common eigenline, and at zero variables the eigenvalue of every $T(t)$ on it is 1 as well (take the empty-empty entries of $T(t)T(\gamma)^m$). 

The linear terms of the common-line eigenvalue vanish because at zero variables the right and dual lines of the projection are in the spin-zero site sector (by total-spin conservation and the nonzero empty-empty entry); linear terms of $T(t)$ change the site's total spin by 1 up or down. Thus the eigenvalue is of the form $1+e_\ell e_r q(t;\xi)$; here $q$ can be computed by contraction of $T(t)$ with vectors independent of $t$ (or trace against the rank-one projection including nilpotents, for instance defined by the Taylor expansion at zero boundary variables of an ordinary simple-pole resolvent projection). At the homogeneous rapidities, the spectral decomposition gives
\[
 [e_\ell e_r]\,(T(\gamma)^m)_{\mathbf0,\mathbf0}
 =m q(\gamma;\xi)+O(1).
\]
Indeed $(1+e_\ell e_rq)^m=1+m e_\ell e_rq$. The mixed coefficient of the continued eigenprojection contributes a term independent of $m$, and derivatives of the remaining spectral part decay by the gap.

There is also a direct path interpretation of this coefficient. Since $l(\gamma)=1$, it sums paths with one positive-arrow terminal on each wall of the $2m$-row rectangle, directed right-to-left and with phase one. Translation averaging and exhaustion therefore give
\[
 B_N=\lim_{m\to\infty}\frac{1}{2m}
 [e_\ell e_r]\,(T(\gamma)^m)_{\mathbf0,\mathbf0}
 =\frac12 q(\gamma;\xi).
\]
The factor two counts the rows in one transfer step.

\paragraph{Rational dependence on the spectral parameter.}
We have identified $q(\gamma)=2B_N$ at the homogeneous physical specialization. Its spectral-parameter dependence now supplies the interpolation conditions. For fixed $\xi$, $q(t;\xi)$ is invariant by $t\mapsto b-t,t+\pi$, by the identities for $T$. From $T$ its denominator thus divides, with $u=t-\gamma$ and $X=\cos(2u)$,
\[
 (X+\cos 2L)\prod_i (X+\cos(2\xi_i+L))(X+\cos(2\xi_i-L)).
\]
Indeed these factors are the paired sine denominators of $K,R$ at reflected arguments, giving polynomials in $X$ after clearing them. As $t=\gamma+iY$, $Y\to\infty$, the coefficient of growth in $q$ comes only from the $e_\ell e_r$ part of $T(t)$ since the $R$'s and $k$'s stay bounded and $l(t),l(b-t)\sim\exp(Y)$. Divided by $l(t)l(b-t)$ this part tends on total site spin 0 to twice identity. In fact the first row's $R$ has $v\to1$, $a,h\to0$, $d\to c^{-1}$ in the sense of reciprocal limiting phases, and has only the possible limiting off-diagonal transition of auxiliary spin $+1$ to $-1$ ($c+d/s^2\to0$); the other row is reversed. The auxiliary spin pairs in this coefficient must both be $+0$ or $0+$, so only diagonal elements survive and their products are 1 on the indicated sector. Hence $q\sim 4 X$.

Putting $p(u)=q(\gamma+u;\xi)$ now at homogeneous $\xi$, we have obtained a rational function of $X$ with polynomial part $4X+\rho$ and denominator dividing $(X+1)^N(X+\cos 2L)^{N+1}$. From $T(0),T(4L)$ we have $p(3L/2)=p(5L/2)=0$. From \eqref{eq:legacy-P-6}, before taking the homogeneous limit, the even function
\[
 q(\gamma+u-L;\xi)+q(\gamma+u+L;\xi)-q(\gamma+u;\xi)
\]
vanishes at $u=L/2\pm\xi_i$. It is analytic near $u=0,L$ for $\xi$ near the homogeneous point. Thus in the homogeneous limit $p(u-L)+p(u+L)-p(u)$ has a zero of multiplicity at least $2N$ at 0 and at least $N$ at $L$ (use generic distinct nodes first, including the even reflections of the ones near 0).

\end{proof}

\subsection{The finite functional}

\begin{proof}[Proof of Theorem~\ref{thm:strip-functional}]
The scalar function constructed in Section~\ref{sec:strip-scalar} can now be represented by the asserted functional. Partial fractions give, on $|\operatorname{Re}u|\le5L/2$, $p(u)=\mathcal L(C_u/D)$ for some $f\in M+\mathbf C g$ and $\rho$ as in \eqref{eq:legacy-P-1}. Indeed
\[
 \int_0^\infty\frac{\cosh(uz)\sinh(mz)}{\sinh(\pi z/2)}\,dz
 =\frac{\tan(m+u)+\tan(m-u)}2=\frac{\sin 2m}{X+\cos 2m},
\]
near $m=0,L$. The integral follows e.g. by geometric expansion of the denominator and the partial fractions for $\tan y$ (from the ordinary sine product or a rectangular residue contour); convergence is absolute. Derivatives $1,3,\ldots,2N-1$ at $m=0$ and $0,1,\ldots,N$ at $L$ span precisely the required polar parts. After division by the factor $z$ in the measure, the first family gives $1,x,\ldots,x^{N-1}$. In the second family, derivative order zero gives $g$, and orders $1,\ldots,N$ give the derivatives of $C_m$ at $m=L$ spanning the other half of $M$. The atoms at $-4$ and $0$ supply $4X$ and $\rho$, respectively.

Then
\[
 p(u-L)+p(u+L)-p(u)=\mathcal L(C_u),
\]
so the multiplicities give $\mathcal L(M)=0$, and $p(3L/2),p(5L/2)=0$ give the other vanishing conditions since $C_{3L/2}/D=C_{L/2}$ and $C_{5L/2}/D=C_{3L/2}+C_{L/2}-C_{L/2}/D$. This completes \eqref{eq:legacy-P-1}.

\end{proof}
